% ID: SEL-D0077-OA-E03
% Tanda S077.
\textbf{Parte I. a)} Hay $3+2=5$ puestos con ordenador averiado entre 20 equiprobables. Por tanto, $P(\text{no funciona})=5/20=\boxed{1/4}$.

\textbf{b)} De los 15 puestos que funcionan, 7 son compartidos y 8 individuales. Condicionando a que el ordenador funcione, la probabilidad pedida es $\boxed{7/15}$.

\textbf{Parte II. a)} Al tomar muestras aleatorias independientes de una población normal, la media muestral también es normal. Su media es 28 kg y su desviación típica es $2{,}7/\sqrt9=0{,}9$ kg. Escribiendo el segundo parámetro como varianza, $\bar X\sim N(28;\ 0{,}9^2)$.

\textbf{b)} Tipificamos con $Z=(\bar X-28)/0{,}9$:
\[
 \begin{aligned}
 P(26\le\bar X\le29)
 &=P\left(-\frac{20}{9}\le Z\le\frac{10}{9}\right)\\
 &=\Phi(10/9)-\Phi(-20/9)\approx\boxed{0{,}8536}.
 \end{aligned}
\]
La probabilidad es aproximadamente del $85{,}36\%$. Al ser una distribución continua, incluir los extremos no cambia el resultado.
